\documentclass[11pt,a4paper]{article}
\usepackage[margin=1in,landscape]{geometry}
\usepackage{booktabs}
\usepackage{pdflscape}
\usepackage{makecell}
\usepackage{placeins}
\usepackage{caption}
\usepackage{amsmath,amssymb}
\usepackage{hyperref}
\title{\textbf{Section 3.1 IM-OLS Cointegration Estimation Matrix: U.S. NFC Transformation Elasticity (1929--2024)}}
\author{\textbf{Department of Economics, UMass Amherst}}
\date{August 2026}
\begin{document}
\maketitle

\section*{Data Provenance and Estimation Summary}
This report presents the empirical cointegration results for Section 3.1 of Chapter 2. All models are estimated via Cointegrating Multivariate Polynomial Regression (CMPR) IM-OLS (Vogelsang et al. 2024) over the full NIPA sample window (1931--2024, $T=94$ observations after listwise deletion for the lag-1 conditioning variables $\omega_{t-1}$ and $\iota_{t-1}$). Point estimates are tuning-parameter free. To ensure exact size control under the non-linear integration orders of the regressors ($I(1) \times I(1)$ and $I(1) \times I(1) \times I(1)$) and account for break-search degrees of freedom loss, significance stars ($^*p<0.10, ^{**}p<0.05, ^{***}p<0.01$) for all coefficient rows ($\hat{\alpha}$, $\hat{\theta}_0$, $\hat{\theta}_1$, $\hat{\beta}_{\text{cm}}$) and the Residual VW Type 2 ($S_{\text{Type 2}}$ / $CT_{\text{IM}}$) cointegration test are computed using a **100,000-pass custom Monte Carlo simulation suite** tailored to each model's exact DGP specification. Distributional conditioning variables ($\omega_{t-1}$ and $\iota_{t-1}$) are centered at their 1929 initial values, preserving $\hat{\theta}_0$ as the capital elasticity at the initial state. VIF statistics are computed on the level regressors space. Step level-shift dummies $D^{\text{step}}_{T_b}$ represent structural breaks identified via a first-pass Bai-Perron residual audit (15\% trimming) and are dynamically adjusted within the custom critical value DGP.

\FloatBarrier
\section{Specification Set 1: Double Interactions ($k^{\text{cap}} \cdot \omega_{t-1}$)}
\FloatBarrier
\begin{table}[H]
\centering
\caption{Specification Set 1 (Part A: Baseline NFC \& Corporate GVA Wage Shares): Double Interaction IM-OLS Cointegration Estimates}
\label{tab:spec_set_1_a}
\scriptsize
\begin{tabular}{p{3.8cm}p{1.5cm}p{1.5cm}p{1.5cm}p{1.5cm}p{1.5cm}p{1.5cm}p{1.5cm}p{1.5cm}p{1.5cm}p{1.5cm}p{1.5cm}p{1.5cm}}
\toprule
\multicolumn{13}{l}{\textbf{Dependent Variable: $y_t^{\text{NFC}}$ (Log NFC Real GVA)}} \\
\midrule
Model & M1.1a & M1.1b & M1.2a & M1.2b & M1.3a & M1.3b & M1.4a & M1.4b & M1.5a & M1.5b & M1.6a & M1.6b  \\
Demand Control & None & None & NetC Share & NetC Share & None & None & NetC Share & NetC Share & None & None & NetC Share & NetC Share  \\
\midrule
$\hat{\alpha}$ (Intercept)  & 6.47$^{*}$ & 8.97 & 6.47$^{*}$ & 8.97 & 6.29$^{*}$ & 6.81 & 6.29$^{*}$ & 6.81 & 6.48$^{*}$ & 9.19 & 6.48$^{*}$ & 9.19  \\
 & {\scriptsize (0.34)} & {\scriptsize (0.89)} & {\scriptsize (0.34)} & {\scriptsize (0.89)} & {\scriptsize (0.34)} & {\scriptsize (0.83)} & {\scriptsize (0.34)} & {\scriptsize (0.83)} & {\scriptsize (0.35)} & {\scriptsize (0.98)} & {\scriptsize (0.35)} & {\scriptsize (0.98)}  \\
$\hat{\theta}_0^{\text{priv}}$ (Private Core $k_{priv}$)  & 0.59$^{**}$ & 0.34 & 0.59$^{**}$ & 0.34 & --- & --- & --- & --- & 0.59$^{**}$ & 0.33 & 0.59$^{**}$ & 0.33  \\
 & {\scriptsize (0.02)} & {\scriptsize (0.08)} & {\scriptsize (0.02)} & {\scriptsize (0.08)} &  &  &  &  & {\scriptsize (0.03)} & {\scriptsize (0.09)} & {\scriptsize (0.03)} & {\scriptsize (0.09)}  \\
$\hat{\theta}_0^{\text{aug}}$ (Infra-Augmented $k_{aug}$)  & --- & --- & --- & --- & 0.59$^{**}$ & 0.52 & 0.59$^{**}$ & 0.52 & --- & --- & --- & ---  \\
 &  &  &  &  & {\scriptsize (0.02)} & {\scriptsize (0.07)} & {\scriptsize (0.02)} & {\scriptsize (0.07)} &  &  &  &   \\
$\hat{\theta}_1$ (Interaction)  & 0.02 & -0.05 & 0.02 & -0.05 & 0.06 & -0.07 & 0.06 & -0.07 & 0.01 & -0.07 & 0.01 & -0.07  \\
 & {\scriptsize (0.13)} & {\scriptsize (0.15)} & {\scriptsize (0.13)} & {\scriptsize (0.15)} & {\scriptsize (0.12)} & {\scriptsize (0.10)} & {\scriptsize (0.12)} & {\scriptsize (0.10)} & {\scriptsize (0.12)} & {\scriptsize (0.16)} & {\scriptsize (0.12)} & {\scriptsize (0.16)}  \\
Step Shifts ($D^{\text{step}}_{T_b}$)  & No & Yes & No & Yes & No & Yes & No & Yes & No & Yes & No & Yes  \\
\midrule
\multicolumn{13}{l}{\underline{Collinearity Diagnostics (VIF)}} \\
VIF ($k^{\text{cap}}$) & 1.38 & 1.38 & 1.38 & 1.38 & 1.40 & 1.40 & 1.40 & 1.40 & 1.54 & 1.54 & 1.54 & 1.54  \\
VIF (Interaction) & 1.38 & 1.38 & 1.38 & 1.38 & 1.40 & 1.40 & 1.40 & 1.40 & 1.54 & 1.54 & 1.54 & 1.54  \\
\midrule
\multicolumn{13}{l}{\underline{Cointegration Test Triad}} \\
\makecell[l]{Resid. ADF ($t_{\text{ADF}}$) \\ {\scriptsize [$H_0$: No Coint.]}} & -4.15$^{**}$ & -3.00 & -4.15$^{**}$ & -3.00 & -3.97$^{**}$ & -3.52$^{*}$ & -3.97$^{**}$ & -3.52$^{*}$ & -4.10$^{**}$ & -3.15 & -4.10$^{**}$ & -3.15  \\
\makecell[l]{Shin LM ($C_{\text{Shin}}$) \\ {\scriptsize [$H_0$: Coint.]}} & 2.12$^{***}$ & 0.56$^{***}$ & 2.12$^{***}$ & 0.56$^{***}$ & 2.15$^{***}$ & 0.92$^{***}$ & 2.15$^{***}$ & 0.92$^{***}$ & 2.08$^{***}$ & 0.45$^{***}$ & 2.08$^{***}$ & 0.45$^{***}$  \\
\makecell[l]{VW Type 2 ($S_{\text{Type 2}}$) \\ {\scriptsize [$H_0$: Coint.]}} & 819.23$^{***}$ & 1835.00$^{***}$ & 819.23$^{***}$ & 1835.00$^{***}$ & 924.93$^{***}$ & 719.49$^{***}$ & 924.93$^{***}$ & 719.49$^{***}$ & 934.98$^{***}$ & 1794.32$^{***}$ & 934.98$^{***}$ & 1794.32$^{***}$  \\
\midrule
Bandwidth Ratio ($b$) & 0.2852 & 0.1737 & 0.2852 & 0.1737 & 0.2832 & 0.1946 & 0.2832 & 0.1946 & 0.2876 & 0.1658 & 0.2876 & 0.1658  \\
Sample Size ($T$) & 94 & 94 & 94 & 94 & 94 & 94 & 94 & 94 & 94 & 94 & 94 & 94  \\
\bottomrule
\multicolumn{13}{p{\linewidth}}{\footnotesize Note: Estimated via Cointegrating Multivariate Polynomial Regression (CMPR) IM-OLS (Vogelsang et al. 2024) over 1929--2024. Standard errors in parentheses. Significance: $^*p<0.10, ^{**}p<0.05, ^{***}p<0.01$. Stars for coefficient rows ($\hat{\alpha}$, $\hat{\theta}_0$, $\hat{\theta}_1$, $\hat{\beta}_{\text{cm}}$) and the Residual VW Type 2 ($S_{\text{Type 2}}$ / $CT_{\text{IM}}$) cointegration test are derived from a **100,000-pass custom Monte Carlo simulation suite** mapping the exact DGP variant (double vs. triple interaction polynomials, with/without step level-shifts, and with/without deterministic $I(0)$ demand controls). This custom simulation adjusts the critical values for the non-linear integration orders of the regressors ($I(1) \times I(1)$ and $I(1) \times I(1) \times I(1)$) and accounts for break-search degrees of freedom loss under the dynamic Bai-Perron audit (15\% trimming). Resid. ADF stars follow MacKinnon (2010) Engle-Granger critical values, indexed by the total number of $I(1)$ variables; the test is lag-augmented with the order selected by AIC (max 8) to whiten serial correlation. Shin LM stars use Shin's (1994) critical values (10\%/5\%/1\% = 0.095/0.121/0.184) and are included for reference. **Star reversal warning:** For Residual ADF, stars indicate rejection of No Cointegration (supporting cointegration). For Shin LM and VW Type 2, stars indicate rejection of Cointegration. Distributional variables centered at 1929 initial state ($t_0$).}
\end{tabular}
\end{table}


\FloatBarrier
\begin{table}[H]
\centering
\caption{Specification Set 1 (Part B: Shaikh-Tonak Tier 2 Wage Share): Double Interaction IM-OLS Cointegration Estimates}
\label{tab:spec_set_1_b}
\scriptsize
\begin{tabular}{p{3.8cm}p{1.5cm}p{1.5cm}p{1.5cm}p{1.5cm}p{1.5cm}p{1.5cm}p{1.5cm}p{1.5cm}p{1.5cm}p{1.5cm}p{1.5cm}p{1.5cm}}
\toprule
\multicolumn{13}{l}{\textbf{Dependent Variable: $y_t^{\text{NFC}}$ (Log NFC Real GVA)}} \\
\midrule
Model & M1.7a & M1.7b & M1.8a & M1.8b & M1.9a & M1.9b & M1.10a & M1.10b & M1.11a & M1.11b & M1.12a & M1.12b  \\
Demand Control & None & None & NetC Share & NetC Share & None & None & NetC Share & NetC Share & None & None & NetC Share & NetC Share  \\
\midrule
$\hat{\alpha}$ (Intercept)  & 6.29$^{*}$ & 8.75 & 6.29$^{*}$ & 8.75 & 7.64 & 5.44 & 7.64 & 5.44 & 7.52 & 4.84 & 7.52 & 4.84  \\
 & {\scriptsize (0.35)} & {\scriptsize (0.90)} & {\scriptsize (0.35)} & {\scriptsize (0.90)} & {\scriptsize (0.67)} & {\scriptsize (1.23)} & {\scriptsize (0.67)} & {\scriptsize (1.23)} & {\scriptsize (0.69)} & {\scriptsize (1.24)} & {\scriptsize (0.69)} & {\scriptsize (1.24)}  \\
$\hat{\theta}_0^{\text{priv}}$ (Private Core $k_{priv}$)  & --- & --- & --- & --- & 0.48 & 0.67 & 0.48 & 0.67 & --- & --- & --- & ---  \\
 &  &  &  &  & {\scriptsize (0.06)} & {\scriptsize (0.11)} & {\scriptsize (0.06)} & {\scriptsize (0.11)} &  &  &  &   \\
$\hat{\theta}_0^{\text{aug}}$ (Infra-Augmented $k_{aug}$)  & 0.59$^{**}$ & 0.35 & 0.59$^{**}$ & 0.35 & --- & --- & --- & --- & 0.48 & 0.70 & 0.48 & 0.70  \\
 & {\scriptsize (0.02)} & {\scriptsize (0.08)} & {\scriptsize (0.02)} & {\scriptsize (0.08)} &  &  &  &  & {\scriptsize (0.06)} & {\scriptsize (0.10)} & {\scriptsize (0.06)} & {\scriptsize (0.10)}  \\
$\hat{\theta}_1$ (Interaction)  & 0.05 & -0.04 & 0.05 & -0.04 & -0.28 & -0.19 & -0.28 & -0.19 & -0.27 & -0.18 & -0.27 & -0.18  \\
 & {\scriptsize (0.12)} & {\scriptsize (0.14)} & {\scriptsize (0.12)} & {\scriptsize (0.14)} & {\scriptsize (0.13)} & {\scriptsize (0.13)} & {\scriptsize (0.13)} & {\scriptsize (0.13)} & {\scriptsize (0.13)} & {\scriptsize (0.13)} & {\scriptsize (0.13)} & {\scriptsize (0.13)}  \\
Step Shifts ($D^{\text{step}}_{T_b}$)  & No & Yes & No & Yes & No & Yes & No & Yes & No & Yes & No & Yes  \\
\midrule
\multicolumn{13}{l}{\underline{Collinearity Diagnostics (VIF)}} \\
VIF ($k^{\text{cap}}$) & 1.56 & 1.56 & 1.56 & 1.56 & 8.04 & 8.04 & 8.04 & 8.04 & 8.14 & 8.14 & 8.14 & 8.14  \\
VIF (Interaction) & 1.56 & 1.56 & 1.56 & 1.56 & 8.04 & 8.04 & 8.04 & 8.04 & 8.14 & 8.14 & 8.14 & 8.14  \\
\midrule
\multicolumn{13}{l}{\underline{Cointegration Test Triad}} \\
\makecell[l]{Resid. ADF ($t_{\text{ADF}}$) \\ {\scriptsize [$H_0$: No Coint.]}} & -3.92$^{**}$ & -2.95 & -3.92$^{**}$ & -2.95 & -3.40 & -2.78 & -3.40 & -2.78 & -3.13 & -2.86 & -3.13 & -2.86  \\
\makecell[l]{Shin LM ($C_{\text{Shin}}$) \\ {\scriptsize [$H_0$: Coint.]}} & 2.11$^{***}$ & 0.60$^{***}$ & 2.11$^{***}$ & 0.60$^{***}$ & 0.76$^{***}$ & 0.20$^{***}$ & 0.76$^{***}$ & 0.20$^{***}$ & 0.91$^{***}$ & 0.20$^{***}$ & 0.91$^{***}$ & 0.20$^{***}$  \\
\makecell[l]{VW Type 2 ($S_{\text{Type 2}}$) \\ {\scriptsize [$H_0$: Coint.]}} & 1004.15$^{***}$ & 2306.76$^{***}$ & 1004.15$^{***}$ & 2306.76$^{***}$ & 743.93$^{***}$ & 2169.68$^{***}$ & 743.93$^{***}$ & 2169.68$^{***}$ & 713.58$^{***}$ & 1404.04$^{***}$ & 713.58$^{***}$ & 1404.04$^{***}$  \\
\midrule
Bandwidth Ratio ($b$) & 0.2856 & 0.1704 & 0.2856 & 0.1704 & 0.2665 & 0.2105 & 0.2665 & 0.2105 & 0.2830 & 0.1907 & 0.2830 & 0.1907  \\
Sample Size ($T$) & 94 & 94 & 94 & 94 & 94 & 94 & 94 & 94 & 94 & 94 & 94 & 94  \\
\bottomrule
\multicolumn{13}{p{\linewidth}}{\footnotesize Note: Estimated via Cointegrating Multivariate Polynomial Regression (CMPR) IM-OLS (Vogelsang et al. 2024) over 1929--2024. Standard errors in parentheses. Significance: $^*p<0.10, ^{**}p<0.05, ^{***}p<0.01$. Stars for coefficient rows ($\hat{\alpha}$, $\hat{\theta}_0$, $\hat{\theta}_1$, $\hat{\beta}_{\text{cm}}$) and the Residual VW Type 2 ($S_{\text{Type 2}}$ / $CT_{\text{IM}}$) cointegration test are derived from a **100,000-pass custom Monte Carlo simulation suite** mapping the exact DGP variant (double vs. triple interaction polynomials, with/without step level-shifts, and with/without deterministic $I(0)$ demand controls). This custom simulation adjusts the critical values for the non-linear integration orders of the regressors ($I(1) \times I(1)$ and $I(1) \times I(1) \times I(1)$) and accounts for break-search degrees of freedom loss under the dynamic Bai-Perron audit (15\% trimming). Resid. ADF stars follow MacKinnon (2010) Engle-Granger critical values, indexed by the total number of $I(1)$ variables; the test is lag-augmented with the order selected by AIC (max 8) to whiten serial correlation. Shin LM stars use Shin's (1994) critical values (10\%/5\%/1\% = 0.095/0.121/0.184) and are included for reference. **Star reversal warning:** For Residual ADF, stars indicate rejection of No Cointegration (supporting cointegration). For Shin LM and VW Type 2, stars indicate rejection of Cointegration. Distributional variables centered at 1929 initial state ($t_0$).}
\end{tabular}
\end{table}


\FloatBarrier
\section{Specification Set 2: Triple Interactions ($k^{\text{cap}} \cdot \omega_{t-1} \cdot \iota_{t-1}$)}
\FloatBarrier
\begin{table}[H]
\centering
\caption{Specification Set 2 (Part A: Baseline NFC \& Corporate GVA Wage Shares): Triple Interaction IM-OLS Cointegration Estimates}
\label{tab:spec_set_2_a}
\scriptsize
\begin{tabular}{p{3.8cm}p{1.5cm}p{1.5cm}p{1.5cm}p{1.5cm}p{1.5cm}p{1.5cm}p{1.5cm}p{1.5cm}p{1.5cm}p{1.5cm}p{1.5cm}p{1.5cm}}
\toprule
\multicolumn{13}{l}{\textbf{Dependent Variable: $y_t^{\text{NFC}}$ (Log NFC Real GVA)}} \\
\midrule
Model & M2.1a & M2.1b & M2.2a & M2.2b & M2.3a & M2.3b & M2.4a & M2.4b & M2.5a & M2.5b & M2.6a & M2.6b  \\
Demand Control & None & None & NetC Share & NetC Share & None & None & NetC Share & NetC Share & None & None & NetC Share & NetC Share  \\
\midrule
$\hat{\alpha}$ (Intercept)  & 6.46 & 8.18 & 9.39 & 9.41 & 6.20 & 8.08 & 9.14 & 7.22 & 6.50 & 7.76 & 9.32 & 9.36  \\
 & {\scriptsize (0.49)} & {\scriptsize (1.06)} & {\scriptsize (0.77)} & {\scriptsize (0.88)} & {\scriptsize (0.50)} & {\scriptsize (1.10)} & {\scriptsize (0.79)} & {\scriptsize (0.88)} & {\scriptsize (0.50)} & {\scriptsize (1.16)} & {\scriptsize (0.75)} & {\scriptsize (0.89)}  \\
$\hat{\theta}_0^{\text{priv}}$ (Private Core $k_{priv}$)  & 0.59 & 0.42 & 0.51 & 0.36 & --- & --- & --- & --- & 0.59 & 0.45 & 0.51 & 0.36  \\
 & {\scriptsize (0.04)} & {\scriptsize (0.09)} & {\scriptsize (0.03)} & {\scriptsize (0.07)} &  &  &  &  & {\scriptsize (0.04)} & {\scriptsize (0.10)} & {\scriptsize (0.03)} & {\scriptsize (0.07)}  \\
$\hat{\theta}_0^{\text{aug}}$ (Infra-Augmented $k_{aug}$)  & --- & --- & --- & --- & 0.60 & 0.41 & 0.52 & 0.50 & --- & --- & --- & ---  \\
 &  &  &  &  & {\scriptsize (0.04)} & {\scriptsize (0.09)} & {\scriptsize (0.03)} & {\scriptsize (0.06)} &  &  &  &   \\
$\hat{\theta}_1$ (Interaction)  & 0.00 & -0.19 & -0.06 & -0.09 & 0.04 & -0.17 & -0.03 & -0.46 & -0.01 & -0.16 & -0.06 & -0.07  \\
 & {\scriptsize (0.10)} & {\scriptsize (0.09)} & {\scriptsize (0.07)} & {\scriptsize (0.08)} & {\scriptsize (0.10)} & {\scriptsize (0.10)} & {\scriptsize (0.07)} & {\scriptsize (0.13)} & {\scriptsize (0.10)} & {\scriptsize (0.10)} & {\scriptsize (0.07)} & {\scriptsize (0.08)}  \\
$\hat{\beta}_{\text{cm}}$ ($\text{NetC}^{\text{share}}$)  & --- & --- & -3.45 & -0.93 & --- & --- & -3.42 & -0.25 & --- & --- & -3.37 & -0.92  \\
 &  &  & {\scriptsize (0.76)} & {\scriptsize (0.62)} &  &  & {\scriptsize (0.77)} & {\scriptsize (0.60)} &  &  & {\scriptsize (0.75)} & {\scriptsize (0.63)}  \\
Step Shifts ($D^{\text{step}}_{T_b}$)  & No & Yes & No & Yes & No & Yes & No & Yes & No & Yes & No & Yes  \\
\midrule
\multicolumn{13}{l}{\underline{Collinearity Diagnostics (VIF)}} \\
VIF ($k^{\text{cap}}$) & 3.57 & 3.57 & 4.41 & 4.41 & 3.69 & 3.69 & 4.61 & 4.61 & 3.86 & 3.86 & 4.64 & 4.64  \\
VIF (Interaction) & 3.57 & 3.57 & 3.58 & 3.58 & 3.69 & 3.69 & 3.71 & 3.71 & 3.86 & 3.86 & 3.87 & 3.87  \\
\midrule
\multicolumn{13}{l}{\underline{Cointegration Test Triad}} \\
\makecell[l]{Resid. ADF ($t_{\text{ADF}}$) \\ {\scriptsize [$H_0$: No Coint.]}} & -4.02$^{**}$ & -4.00$^{**}$ & -2.11 & -2.64 & -3.82$^{**}$ & -3.00 & -2.16 & -3.21 & -3.97$^{**}$ & -3.93$^{**}$ & -2.10 & -2.62  \\
\makecell[l]{Shin LM ($C_{\text{Shin}}$) \\ {\scriptsize [$H_0$: Coint.]}} & 1.86$^{***}$ & 0.66$^{***}$ & 2.08$^{***}$ & 0.83$^{***}$ & 1.90$^{***}$ & 0.71$^{***}$ & 2.08$^{***}$ & 0.41$^{***}$ & 1.78$^{***}$ & 0.72$^{***}$ & 2.08$^{***}$ & 0.83$^{***}$  \\
\makecell[l]{VW Type 2 ($S_{\text{Type 2}}$) \\ {\scriptsize [$H_0$: Coint.]}} & 3910.69$^{***}$ & 276.16$^{***}$ & --- & --- & 3724.48$^{***}$ & 359.15$^{***}$ & --- & --- & 4021.16$^{***}$ & 140.50$^{***}$ & --- & ---  \\
\midrule
Bandwidth Ratio ($b$) & 0.2908 & 0.1688 & NA & NA & 0.2916 & 0.1688 & NA & NA & 0.2934 & 0.1734 & NA & NA  \\
Sample Size ($T$) & 94 & 94 & 94 & 94 & 94 & 94 & 94 & 94 & 94 & 94 & 94 & 94  \\
\bottomrule
\multicolumn{13}{p{\linewidth}}{\footnotesize Note: Estimated via Cointegrating Multivariate Polynomial Regression (CMPR) IM-OLS (Vogelsang et al. 2024) over 1929--2024. Standard errors in parentheses. Significance: $^*p<0.10, ^{**}p<0.05, ^{***}p<0.01$. Stars for coefficient rows ($\hat{\alpha}$, $\hat{\theta}_0$, $\hat{\theta}_1$, $\hat{\beta}_{\text{cm}}$) and the Residual VW Type 2 ($S_{\text{Type 2}}$ / $CT_{\text{IM}}$) cointegration test are derived from a **100,000-pass custom Monte Carlo simulation suite** mapping the exact DGP variant (double vs. triple interaction polynomials, with/without step level-shifts, and with/without deterministic $I(0)$ demand controls). This custom simulation adjusts the critical values for the non-linear integration orders of the regressors ($I(1) \times I(1)$ and $I(1) \times I(1) \times I(1)$) and accounts for break-search degrees of freedom loss under the dynamic Bai-Perron audit (15\% trimming). Resid. ADF stars follow MacKinnon (2010) Engle-Granger critical values, indexed by the total number of $I(1)$ variables; the test is lag-augmented with the order selected by AIC (max 8) to whiten serial correlation. Shin LM stars use Shin's (1994) critical values (10\%/5\%/1\% = 0.095/0.121/0.184) and are included for reference. **Star reversal warning:** For Residual ADF, stars indicate rejection of No Cointegration (supporting cointegration). For Shin LM and VW Type 2, stars indicate rejection of Cointegration. Distributional variables centered at 1929 initial state ($t_0$).}
\end{tabular}
\end{table}


\FloatBarrier
\begin{table}[H]
\centering
\caption{Specification Set 2 (Part B: Shaikh-Tonak Tier 2 Wage Share): Triple Interaction IM-OLS Cointegration Estimates}
\label{tab:spec_set_2_b}
\scriptsize
\begin{tabular}{p{3.8cm}p{1.5cm}p{1.5cm}p{1.5cm}p{1.5cm}p{1.5cm}p{1.5cm}p{1.5cm}p{1.5cm}p{1.5cm}p{1.5cm}p{1.5cm}p{1.5cm}}
\toprule
\multicolumn{13}{l}{\textbf{Dependent Variable: $y_t^{\text{NFC}}$ (Log NFC Real GVA)}} \\
\midrule
Model & M2.7a & M2.7b & M2.8a & M2.8b & M2.9a & M2.9b & M2.10a & M2.10b & M2.11a & M2.11b & M2.12a & M2.12b  \\
Demand Control & None & None & NetC Share & NetC Share & None & None & NetC Share & NetC Share & None & None & NetC Share & NetC Share  \\
\midrule
$\hat{\alpha}$ (Intercept)  & 6.24 & 7.93 & 9.06 & 8.60 & 8.30 & 8.68 & 9.42 & 9.49 & 8.15 & 8.44 & 9.38 & 8.98  \\
 & {\scriptsize (0.52)} & {\scriptsize (1.13)} & {\scriptsize (0.78)} & {\scriptsize (1.00)} & {\scriptsize (0.81)} & {\scriptsize (0.70)} & {\scriptsize (0.70)} & {\scriptsize (0.81)} & {\scriptsize (0.87)} & {\scriptsize (0.73)} & {\scriptsize (0.74)} & {\scriptsize (0.88)}  \\
$\hat{\theta}_0^{\text{priv}}$ (Private Core $k_{priv}$)  & --- & --- & --- & --- & 0.42 & 0.39 & 0.44 & 0.43 & --- & --- & --- & ---  \\
 &  &  &  &  & {\scriptsize (0.07)} & {\scriptsize (0.06)} & {\scriptsize (0.04)} & {\scriptsize (0.06)} &  &  &  &   \\
$\hat{\theta}_0^{\text{aug}}$ (Infra-Augmented $k_{aug}$)  & 0.60 & 0.43 & 0.52 & 0.41 & --- & --- & --- & --- & 0.43 & 0.40 & 0.44 & 0.45  \\
 & {\scriptsize (0.04)} & {\scriptsize (0.10)} & {\scriptsize (0.03)} & {\scriptsize (0.08)} &  &  &  &  & {\scriptsize (0.07)} & {\scriptsize (0.06)} & {\scriptsize (0.05)} & {\scriptsize (0.07)}  \\
$\hat{\theta}_1$ (Interaction)  & 0.02 & -0.15 & -0.03 & -0.01 & -0.31 & -0.38 & -0.23 & -0.30 & -0.30 & -0.36 & -0.20 & -0.29  \\
 & {\scriptsize (0.10)} & {\scriptsize (0.10)} & {\scriptsize (0.07)} & {\scriptsize (0.09)} & {\scriptsize (0.13)} & {\scriptsize (0.11)} & {\scriptsize (0.08)} & {\scriptsize (0.07)} & {\scriptsize (0.13)} & {\scriptsize (0.11)} & {\scriptsize (0.08)} & {\scriptsize (0.07)}  \\
$\hat{\beta}_{\text{cm}}$ ($\text{NetC}^{\text{share}}$)  & --- & --- & -3.33 & -0.89 & --- & --- & -2.21 & -2.12 & --- & --- & -2.48 & -1.99  \\
 &  &  & {\scriptsize (0.76)} & {\scriptsize (0.63)} &  &  & {\scriptsize (0.66)} & {\scriptsize (0.39)} &  &  & {\scriptsize (0.68)} & {\scriptsize (0.40)}  \\
Step Shifts ($D^{\text{step}}_{T_b}$)  & No & Yes & No & Yes & No & Yes & No & Yes & No & Yes & No & Yes  \\
\midrule
\multicolumn{13}{l}{\underline{Collinearity Diagnostics (VIF)}} \\
VIF ($k^{\text{cap}}$) & 4.00 & 4.00 & 4.84 & 4.84 & 14.01 & 14.01 & 14.05 & 14.05 & 14.46 & 14.46 & 14.48 & 14.48  \\
VIF (Interaction) & 4.00 & 4.00 & 4.01 & 4.01 & 14.01 & 14.01 & 15.20 & 15.20 & 14.46 & 14.46 & 15.56 & 15.56  \\
\midrule
\multicolumn{13}{l}{\underline{Cointegration Test Triad}} \\
\makecell[l]{Resid. ADF ($t_{\text{ADF}}$) \\ {\scriptsize [$H_0$: No Coint.]}} & -3.77$^{**}$ & -2.98 & -2.16 & -2.50 & -3.27 & -2.90 & -2.52 & -3.06 & -2.30 & -2.85 & -2.47 & -3.09  \\
\makecell[l]{Shin LM ($C_{\text{Shin}}$) \\ {\scriptsize [$H_0$: Coint.]}} & 1.81$^{***}$ & 0.76$^{***}$ & 2.10$^{***}$ & 0.90$^{***}$ & 0.96$^{***}$ & 0.51$^{***}$ & 1.29$^{***}$ & 0.38$^{***}$ & 1.25$^{***}$ & 0.52$^{***}$ & 1.31$^{***}$ & 0.39$^{***}$  \\
\makecell[l]{VW Type 2 ($S_{\text{Type 2}}$) \\ {\scriptsize [$H_0$: Coint.]}} & 3822.34$^{***}$ & 223.87$^{***}$ & --- & --- & 88.61$^{**}$ & 3570.21$^{***}$ & --- & --- & 91.10$^{**}$ & 3285.32$^{***}$ & --- & ---  \\
\midrule
Bandwidth Ratio ($b$) & 0.2940 & 0.1691 & NA & NA & 0.2713 & 0.2047 & NA & NA & 0.2924 & 0.2029 & NA & NA  \\
Sample Size ($T$) & 94 & 94 & 94 & 94 & 94 & 94 & 94 & 94 & 94 & 94 & 94 & 94  \\
\bottomrule
\multicolumn{13}{p{\linewidth}}{\footnotesize Note: Estimated via Cointegrating Multivariate Polynomial Regression (CMPR) IM-OLS (Vogelsang et al. 2024) over 1929--2024. Standard errors in parentheses. Significance: $^*p<0.10, ^{**}p<0.05, ^{***}p<0.01$. Stars for coefficient rows ($\hat{\alpha}$, $\hat{\theta}_0$, $\hat{\theta}_1$, $\hat{\beta}_{\text{cm}}$) and the Residual VW Type 2 ($S_{\text{Type 2}}$ / $CT_{\text{IM}}$) cointegration test are derived from a **100,000-pass custom Monte Carlo simulation suite** mapping the exact DGP variant (double vs. triple interaction polynomials, with/without step level-shifts, and with/without deterministic $I(0)$ demand controls). This custom simulation adjusts the critical values for the non-linear integration orders of the regressors ($I(1) \times I(1)$ and $I(1) \times I(1) \times I(1)$) and accounts for break-search degrees of freedom loss under the dynamic Bai-Perron audit (15\% trimming). Resid. ADF stars follow MacKinnon (2010) Engle-Granger critical values, indexed by the total number of $I(1)$ variables; the test is lag-augmented with the order selected by AIC (max 8) to whiten serial correlation. Shin LM stars use Shin's (1994) critical values (10\%/5\%/1\% = 0.095/0.121/0.184) and are included for reference. **Star reversal warning:** For Residual ADF, stars indicate rejection of No Cointegration (supporting cointegration). For Shin LM and VW Type 2, stars indicate rejection of Cointegration. Distributional variables centered at 1929 initial state ($t_0$).}
\end{tabular}
\end{table}


\end{document}
